\chapter{Banach空間}

\paragraph{本章の目標}
ここまでで，Lebesgue積分によって関数を積分し，
収束定理によって極限と積分の交換も扱えるようになった．
次の段階では，関数そのものを「点」とみなし，
関数全体の集まりを1つの空間として考える．
その基本概念がノルム空間とBanach空間である．

本章では
\begin{enumerate}
    \item ノルム空間・Banach空間の定義
    \item 有限次元空間 $\RR^d$ の例
    \item 連続関数空間 $C(X)$ が sup-norm で Banach 空間になること
    \item 測度空間上の $L^p(X,m)$ の定義
    \item Hölder の不等式，Minkowski の不等式，Riesz--Fischer の定理
\end{enumerate}
を証明する．

\section{動機づけ}

微分積分では，1つ1つの関数を個別に扱ってきた．
しかし関数解析では，
\[
\text{連続関数全体},\qquad
\text{可積分関数全体},\qquad
\text{二乗可積分関数全体}
\]
のような「関数の集まり」自体を対象にする．

そのとき最初に必要になるのは，
関数の大きさや2つの関数の近さを測る道具である．
例えば連続関数 $f,g$ に対して
\[
\sup_{x\in X}|f(x)-g(x)|
\]
が小さければ，$f$ と $g$ はどの点でも一様に近いと言える．
また可積分関数に対して
\[
\left(\int_X |f-g|^p\,dm\right)^{1/p}
\]
が小さければ，平均的な意味で近いと言える．

この「大きさ」を抽象化したものがノルムであり，
そのノルムに関して極限が取り尽くされている空間がBanach空間である．
Lebesgue積分論の収束定理は，
まさにこうした完備な関数空間の上で本領を発揮する．

\section{ノルム空間とBanach空間}

\begin{definition}[ノルム]
実ベクトル空間 $X$ 上の写像
\[
\|\cdot\|:X\to [0,\infty)
\]
がノルムであるとは，任意の $x,y\in X$ と任意の $a\in\RR$ に対して
次を満たすことをいう．
\begin{enumerate}
    \item
    \[
    \|x\|=0 \iff x=0
    \]
    \item
    \[
    \|ax\|=|a|\,\|x\|
    \]
    \item
    \[
    \|x+y\|\le \|x\|+\|y\|
    \]
\end{enumerate}
\end{definition}

\begin{definition}[ノルム空間]
実ベクトル空間 $X$ とその上のノルム $\|\cdot\|$ の組
\[
(X,\|\cdot\|)
\]
をノルム空間という．
\end{definition}

\begin{remark}
ノルムが与えられると，
\[
d(x,y):=\|x-y\|
\]
によって距離が定まる．
したがってノルム空間では，点列の収束やCauchy列を距離空間と同様に定義できる．
\end{remark}

\begin{definition}[Banach空間]
ノルム空間 $X$ がBanach空間であるとは，
$X$ の任意のCauchy列が $X$ の元に収束することをいう．
\end{definition}

\begin{example}[\texorpdfstring{$\RR^d$}{Rd} の $p$-ノルム]
$x=(x_1,\dots,x_d)\in\RR^d$ に対し，
$1\le p<\infty$ のとき
\[
\|x\|_p:=\left(\sum_{k=1}^d |x_k|^p\right)^{1/p},
\]
また
\[
\|x\|_\infty:=\max_{1\le k\le d}|x_k|
\]
と定める．
これらは $\RR^d$ 上のノルムである．
\end{example}

\begin{proposition}
$1\le p\le\infty$ に対して，
\[
(\RR^d,\|\cdot\|_p)
\]
はBanach空間である．
\end{proposition}
\begin{proof}
$\{x^{(n)}\}_{n=1}^{\infty}$ を $\RR^d$ のCauchy列とし，
\[
x^{(n)}=(x_1^{(n)},\dots,x_d^{(n)})
\]
と書く．
各成分 $k=1,\dots,d$ に対して
\[
|x_k^{(n)}-x_k^{(m)}|
\le
\|x^{(n)}-x^{(m)}\|_p
\]
が成り立つ．
したがって各実数列 $\{x_k^{(n)}\}_{n=1}^{\infty}$ は $\RR$ でCauchy列であり，
$\RR$ の完備性より極限
\[
x_k:=\lim_{n\to\infty}x_k^{(n)}
\]
が存在する．

\[
x:=(x_1,\dots,x_d)\in\RR^d
\]
とおく．
$1\le p<\infty$ のとき
\[
\|x^{(n)}-x\|_p^p
=
\sum_{k=1}^d |x_k^{(n)}-x_k|^p
\to 0
\]
である．
実際，和は有限個しかないので，各項が $0$ に収束すれば和も $0$ に収束する．

また $p=\infty$ のときは，各 $k$ について
\[
|x_k^{(n)}-x_k|\to 0
\]
であるから，有限個の実数の最大値も $0$ に収束して
\[
\|x^{(n)}-x\|_\infty
=
\max_{1\le k\le d}|x_k^{(n)}-x_k|
\to 0
\]
を得る．

よって $\{x^{(n)}\}$ は $x\in\RR^d$ に収束する．
したがって $\RR^d$ は完備である．
\end{proof}

\begin{remark}
有限次元空間では，各成分ごとに議論できるため，
完備性の証明は比較的素朴である．
しかし関数空間では成分が有限個ではないので，
完備性は自明ではない．
この違いが有限次元と無限次元の大きな分かれ目である．
\end{remark}

\section{\texorpdfstring{$C(X)$}{C(X)} と sup-norm}

以後，$X$ をコンパクト距離空間とする．

\begin{definition}[連続関数空間]
\[
C(X):=\{f:X\to\RR \mid f \text{ は連続}\}
\]
とおく．
\end{definition}

\begin{definition}[sup-norm]
$f\in C(X)$ に対して
\[
\|f\|_\infty:=\sup_{x\in X}|f(x)|
\]
と定める．
\end{definition}

\begin{remark}
$X$ がコンパクトで $f$ が連続なら，最大値の定理より $|f|$ は $X$ 上で最大値をとる．
したがって
\[
\|f\|_\infty<\infty
\]
である．
$X$ のコンパクト性は，sup-norm を考えるうえで重要である．
\end{remark}

\begin{proposition}
\[
\|\cdot\|_\infty
\]
は $C(X)$ 上のノルムである．
\end{proposition}
\begin{proof}
まず
\[
\|f\|_\infty\ge 0
\]
であり，
\[
\|f\|_\infty=0
\]
なら
\[
|f(x)|\le \|f\|_\infty=0
\]
が任意の $x\in X$ で成り立つから $f=0$ である．

次に任意の $a\in\RR$ に対して
\[
\|af\|_\infty
=
\sup_{x\in X}|a||f(x)|
=
|a|\,\|f\|_\infty
\]
である．

最後に任意の $x\in X$ に対して
\[
|f(x)+g(x)|\le |f(x)|+|g(x)|
\]
だから，上限をとれば
\[
\|f+g\|_\infty\le \|f\|_\infty+\|g\|_\infty
\]
を得る．
\end{proof}

\begin{lemma}[一様極限の連続性]
$f_n\in C(X)$ が
\[
\|f_n-f\|_\infty\to 0
\]
を満たすとする．
このとき $f$ は連続である．
\end{lemma}
\begin{proof}
$x_0\in X$ を任意に取る．
$\eps>0$ を与える．
$\|f_n-f\|_\infty\to 0$ だから，ある $N$ が存在して
\[
\|f_N-f\|_\infty<\frac{\eps}{3}
\]
である．

$f_N$ は連続だから，ある $\delta>0$ が存在して
\[
d(x,x_0)<\delta
\implies
|f_N(x)-f_N(x_0)|<\frac{\eps}{3}
\]
が成り立つ．
したがって $d(x,x_0)<\delta$ なら
\begin{align*}
|f(x)-f(x_0)|
&\le |f(x)-f_N(x)|+|f_N(x)-f_N(x_0)|+|f_N(x_0)-f(x_0)|\\
&< \frac{\eps}{3}+\frac{\eps}{3}+\frac{\eps}{3}
=\eps.
\end{align*}
よって $f$ は $x_0$ で連続である．
$x_0$ は任意だから，$f$ は $X$ 上連続である．
\end{proof}

\begin{theorem}
\[
(C(X),\|\cdot\|_\infty)
\]
はBanach空間である．
\end{theorem}
\begin{proof}
$\{f_n\}_{n=1}^{\infty}$ を $C(X)$ のCauchy列とする．
すると任意の $\eps>0$ に対して，ある $N$ が存在して
\[
m,n\ge N \implies \|f_n-f_m\|_\infty<\eps
\]
である．

各 $x\in X$ を固定すると
\[
|f_n(x)-f_m(x)|\le \|f_n-f_m\|_\infty
\]
だから，実数列 $\{f_n(x)\}$ はCauchy列である．
ゆえに極限
\[
f(x):=\lim_{n\to\infty}f_n(x)
\]
が存在する．
こうして写像 $f:X\to\RR$ が定まる．

次に $f_n\to f$ が sup-norm で成り立つことを示す．
$\eps>0$ を与え，上のCauchy条件を満たす $N$ を取る．
$n\ge N$ とし，任意の $x\in X$ に対して $m\to\infty$ とすると
\[
|f_n(x)-f(x)|
=
\lim_{m\to\infty}|f_n(x)-f_m(x)|
\le \eps
\]
を得る．
したがって
\[
\|f_n-f\|_\infty\le \eps
\qquad (n\ge N)
\]
であり，
\[
\|f_n-f\|_\infty\to 0
\]
が従う．

最後に補題より，$f$ は連続である．
よって $f\in C(X)$ であり，$\{f_n\}$ は $C(X)$ の元に収束する．
したがって $C(X)$ は完備である．
\end{proof}

\begin{example}
$X=[a,b]$ とすると，
\[
C([a,b])
\]
は sup-norm に関してBanach空間である．
これは微分方程式や積分方程式で最もよく現れる関数空間の1つである．
\end{example}

\section{\texorpdfstring{$L^p(X,m)$}{Lp(X,m)} 空間}

以後，
\[
(X,\calF,m)
\]
を測度空間とし，
$1\le p<\infty$ を固定する．

\begin{definition}
可測関数 $f:X\to\RR$ で
\[
\int_X |f|^p\,dm<\infty
\]
を満たすもの全体を
\[
\mathcal L^p(X,m)
\]
と書く．
\end{definition}

\begin{remark}
$f$ と $g$ が a.e.\ で一致するなら，
第8章と第9章の a.e.\ 不変性より
\[
\int_X |f|^p\,dm=\int_X |g|^p\,dm
\]
である．
したがって $L^p$ では，関数を点ごとにではなく
「a.e.\ の違いを無視した同値類」として扱うのが自然である．
\end{remark}

\begin{definition}[\texorpdfstring{$L^p(X,m)$}{Lp(X,m)}]
$\mathcal L^p(X,m)$ の上に
\[
f\sim g
\iff
f=g \quad \text{a.e.\ on } X
\]
という同値関係を入れる．
この同値類全体を
\[
L^p(X,m)
\]
と書く．
\end{definition}

\begin{definition}[\texorpdfstring{$L^p$}{Lp} ノルム]
$f\in L^p(X,m)$ に対して
\[
\|f\|_p
:=
\left(\int_X |f|^p\,dm\right)^{1/p}
\]
と定める．
\end{definition}

\begin{remark}
厳密には $f$ は同値類であるが，
記号が煩雑になるのを避けるため，
代表元も同じ記号 $f$ で書く．
上の量が代表元の取り方によらないことは a.e.\ 不変性から従う．
\end{remark}

\begin{proposition}
$L^p(X,m)$ は実ベクトル空間である．
\end{proposition}
\begin{proof}
$f,g\in L^p(X,m)$ と $a,b\in\RR$ を取る．
可測性は第7章で示した和とスカラー倍の可測性から従う．

また任意の実数 $u,v$ に対して
\[
|u+v|^p\le (|u|+|v|)^p\le 2^{p-1}(|u|^p+|v|^p)
\]
が成り立つから，
\[
|af+bg|^p
\le
2^{p-1}\bigl(|a|^p|f|^p+|b|^p|g|^p\bigr)
\]
である．
右辺は可積分であるから
\[
\int_X |af+bg|^p\,dm<\infty
\]
を得る．
よって $af+bg\in \mathcal L^p(X,m)$ である．

さらに $f\sim f'$，$g\sim g'$ なら
\[
af+bg=af'+bg'
\quad \text{a.e.}
\]
だから，演算は同値類上でもwell-definedである．
したがって $L^p(X,m)$ は実ベクトル空間である．
\end{proof}

\begin{example}
$X=(0,1]$，$m$ をLebesgue測度とし，
\[
f(x):=x^{-1/2}
\]
とおく．
すると
\[
\int_0^1 |f(x)|\,dx=\int_0^1 x^{-1/2}\,dx=2<\infty
\]
だから
\[
f\in L^1((0,1],m)
\]
である．
しかし
\[
\int_0^1 |f(x)|^2\,dx=\int_0^1 x^{-1}\,dx=\infty
\]
だから
\[
f\notin L^2((0,1],m)
\]
である．
このように，$p$ が変わると空間も変わる．
\end{example}

\section{Hölderの不等式とMinkowskiの不等式}

\begin{lemma}[Youngの不等式]
$1<p<\infty$ とし，
\[
\frac{1}{p}+\frac{1}{q}=1
\]
を満たす $q$ を共役指数とする．
このとき任意の $a,b\ge 0$ に対して
\[
ab\le \frac{a^p}{p}+\frac{b^q}{q}
\]
が成り立つ．
\end{lemma}
\begin{proof}
$a=0$ のときは自明であるから $a>0$ とする．
\[
\phi(t):=at-\frac{t^q}{q}
\qquad (t\ge 0)
\]
とおくと
\[
\phi'(t)=a-t^{q-1}
\]
である．
したがって $\phi$ は
\[
t=a^{1/(q-1)}=a^{p-1}
\]
で最大値をとる．
そこで
\begin{align*}
\phi(t)
&\le \phi(a^{p-1}) \\
&= a\cdot a^{p-1}-\frac{(a^{p-1})^q}{q} \\
&= a^p-\frac{a^p}{q}
=\frac{a^p}{p}.
\end{align*}
$t=b$ とおけば
\[
ab-\frac{b^q}{q}\le \frac{a^p}{p}
\]
すなわち
\[
ab\le \frac{a^p}{p}+\frac{b^q}{q}
\]
を得る．
\end{proof}

\begin{theorem}[Hölderの不等式]
$1<p<\infty$ とし，
$q$ を共役指数とする．
このとき
$f\in L^p(X,m)$，
$g\in L^q(X,m)$ に対して
\[
fg\in L^1(X,m)
\]
であり，
\[
\int_X |fg|\,dm\le \|f\|_p\|g\|_q
\]
が成り立つ．
\end{theorem}
\begin{proof}
$\|f\|_p=0$ または $\|g\|_q=0$ のときは右辺が $0$ であり，
第9章の「積分が $0$ であることの判定」より
$f=0$ a.e.\ または $g=0$ a.e.\ だから自明である．

以下では
\[
\|f\|_p>0,\qquad \|g\|_q>0
\]
とする．
Young の不等式を
\[
a=\frac{|f(x)|}{\|f\|_p},
\qquad
b=\frac{|g(x)|}{\|g\|_q}
\]
に適用すると，a.e.\ で
\[
\frac{|f(x)g(x)|}{\|f\|_p\|g\|_q}
\le
\frac{|f(x)|^p}{p\|f\|_p^p}
+\frac{|g(x)|^q}{q\|g\|_q^q}
\]
を得る．
両辺を積分すると
\begin{align*}
\frac{1}{\|f\|_p\|g\|_q}\int_X |fg|\,dm
&\le
\frac{1}{p\|f\|_p^p}\int_X |f|^p\,dm
+
\frac{1}{q\|g\|_q^q}\int_X |g|^q\,dm \\
&=
\frac{1}{p}+\frac{1}{q}
=1.
\end{align*}
したがって
\[
\int_X |fg|\,dm\le \|f\|_p\|g\|_q.
\]
特に左辺は有限だから $fg\in L^1(X,m)$ である．
\end{proof}

\begin{theorem}[Minkowskiの不等式]
$1\le p<\infty$ とする．
$f,g\in L^p(X,m)$ ならば
\[
f+g\in L^p(X,m)
\]
であり，
\[
\|f+g\|_p\le \|f\|_p+\|g\|_p
\]
が成り立つ．
\end{theorem}
\begin{proof}
まず $p=1$ のときは
\[
|f+g|\le |f|+|g|
\]
を積分すれば
\[
\|f+g\|_1
=
\int_X |f+g|\,dm
\le
\int_X |f|\,dm+\int_X |g|\,dm
=
\|f\|_1+\|g\|_1
\]
である．

以下では $1<p<\infty$ とする．
$q$ を共役指数とする．
$\|f+g\|_p=0$ のときは自明だから，
\[
\|f+g\|_p>0
\]
と仮定する．
すると
\begin{align*}
\|f+g\|_p^p
&=
\int_X |f+g|^p\,dm \\
&=
\int_X |f+g|\,|f+g|^{p-1}\,dm \\
&\le
\int_X |f|\,|f+g|^{p-1}\,dm
+
\int_X |g|\,|f+g|^{p-1}\,dm.
\end{align*}

ここで
\[
\bigl(|f+g|^{p-1}\bigr)^q=|f+g|^p
\]
であるから，Hölder の不等式より
\begin{align*}
\int_X |f|\,|f+g|^{p-1}\,dm
&\le
\|f\|_p
\left(
\int_X |f+g|^p\,dm
\right)^{1/q} \\
&=
\|f\|_p \|f+g\|_p^{p-1},
\end{align*}
同様に
\[
\int_X |g|\,|f+g|^{p-1}\,dm
\le
\|g\|_p \|f+g\|_p^{p-1}.
\]
したがって
\[
\|f+g\|_p^p
\le
(\|f\|_p+\|g\|_p)\|f+g\|_p^{p-1}.
\]
$\|f+g\|_p>0$ だから両辺を $\|f+g\|_p^{p-1}$ で割って
\[
\|f+g\|_p\le \|f\|_p+\|g\|_p
\]
を得る．
\end{proof}

\begin{proposition}
$1\le p<\infty$ に対して
\[
\|\cdot\|_p
\]
は $L^p(X,m)$ 上のノルムである．
\end{proposition}
\begin{proof}
非負性は自明である．
また
\[
\|af\|_p^p
=
\int_X |a|^p|f|^p\,dm
=
|a|^p\|f\|_p^p
\]
だから
\[
\|af\|_p=|a|\,\|f\|_p
\]
である．

さらに $\|f\|_p=0$ なら
\[
\int_X |f|^p\,dm=0
\]
であり，第9章の「積分が $0$ であることの判定」より
\[
|f|^p=0 \quad \text{a.e.}
\]
だから $f=0$ a.e.\ である．
すなわち $L^p$ の元として $f=0$ である．
逆向きは明らかである．

最後に三角不等式は Minkowski の不等式である．
したがって $\|\cdot\|_p$ はノルムである．
\end{proof}

\section{Riesz--Fischerの定理}

\begin{theorem}[Riesz--Fischer]
$1\le p<\infty$ とする．
このとき
\[
(L^p(X,m),\|\cdot\|_p)
\]
はBanach空間である．
\end{theorem}
\begin{proof}
$\{f_n\}_{n=1}^{\infty}$ を $L^p(X,m)$ のCauchy列とする．
各 $f_n$ の代表元を同じ記号で表す．

$\{f_n\}$ はCauchy列だから，
部分列 $\{f_{n_k}\}$ を
\[
\|f_{n_{k+1}}-f_{n_k}\|_p<2^{-k}
\qquad (k\in\NN)
\]
となるように取れる．
各 $k$ に対して
\[
g_k:=|f_{n_{k+1}}-f_{n_k}|
\]
とおく．
すると $g_k\in L^p(X,m)$ である．

さらに
\[
s_N:=\sum_{k=1}^N g_k
\]
とおくと，Minkowski の不等式より
\[
\|s_N\|_p
\le
\sum_{k=1}^N \|g_k\|_p
<
\sum_{k=1}^{\infty}2^{-k}
=1
\]
である．
したがって
\[
\int_X s_N^p\,dm=\|s_N\|_p^p\le 1
\]
である．

$s_N$ は単調増加する非負可測関数列であるから
\[
s(x):=\sum_{k=1}^{\infty}g_k(x)
\]
とおくと，
単調収束定理より
\[
\int_X s^p\,dm
=
\lim_{N\to\infty}\int_X s_N^p\,dm
\le 1.
\]
よって
\[
s\in L^p(X,m)
\]
であり，特に
\[
s(x)<\infty
\quad \text{a.e. on } X
\]
である．

この零集合を除いた各点 $x$ に対して，
任意の $m>\ell$ について
\[
|f_{n_m}(x)-f_{n_\ell}(x)|
\le
\sum_{k=\ell}^{m-1}g_k(x)
\]
が成り立つ．
右辺は $\ell\to\infty$ で $0$ に収束するから，
$\{f_{n_k}(x)\}$ は a.e.\ でCauchy列である．
したがってある可測関数 $f$ が存在して
\[
f_{n_k}(x)\to f(x)
\quad \text{a.e. on } X
\]
となる．

次にこの $f$ が $L^p$ 極限であることを示す．
各 $\ell$ に対して
\[
t_\ell:=\sum_{k=\ell}^{\infty}g_k
=
s-\sum_{k=1}^{\ell-1}g_k
\]
とおくと，
\[
|f-f_{n_\ell}|\le t_\ell
\quad \text{a.e.}
\]
である．
また
\[
\|t_\ell\|_p
\le
\sum_{k=\ell}^{\infty}\|g_k\|_p
\le
\sum_{k=\ell}^{\infty}2^{-k}
\to 0
\]
が Minkowski の不等式から従う．
したがって単調性より
\[
\|f-f_{n_\ell}\|_p\le \|t_\ell\|_p\to 0.
\]
特に $f\in L^p(X,m)$ であり，
\[
f_{n_\ell}\to f \quad \text{in } L^p
\]
が成り立つ．

最後に元の列全体が $f$ に収束することを示す．
$\eps>0$ を与える．
$\{f_n\}$ はCauchy列だから，ある $N_1$ が存在して
\[
n,m\ge N_1
\implies
\|f_n-f_m\|_p<\frac{\eps}{2}
\]
である．
また部分列 $f_{n_k}\to f$ だから，ある $k$ が存在して
\[
n_k\ge N_1,
\qquad
\|f_{n_k}-f\|_p<\frac{\eps}{2}
\]
となる．
すると $n\ge N_1$ に対して
\[
\|f_n-f\|_p
\le
\|f_n-f_{n_k}\|_p+\|f_{n_k}-f\|_p
<
\frac{\eps}{2}+\frac{\eps}{2}
=\eps.
\]
よって
\[
f_n\to f \quad \text{in } L^p(X,m)
\]
である．
したがって $L^p(X,m)$ は完備である．
\end{proof}

\begin{remark}
この証明の要点は，
Cauchy列から「差の和が可積分になる部分列」を取り出し，
その差の級数
\[
\sum_{k=1}^{\infty}|f_{n_{k+1}}-f_{n_k}|
\]
を使って点ごとの極限を構成することである．
これは収束定理と不等式を組み合わせた，
Lebesgue積分論らしい完備性の証明である．
\end{remark}

\section{解釈とまとめ}

Banach空間の本質は，
「極限を取っても空間の外へ出ない」
という点にある．
有限次元の $\RR^d$ ではこれは座標ごとに従ったが，
関数空間では積分論が必要になった．

本章で得た代表的なBanach空間は次の2つである．
\begin{enumerate}
    \item コンパクト空間上の連続関数空間
    \[
    C(X),\qquad \|f\|_\infty=\sup_{x\in X}|f(x)|
    \]
    \item 測度空間上の可積分関数空間
    \[
    L^p(X,m),\qquad
    \|f\|_p=\left(\int_X |f|^p\,dm\right)^{1/p}
    \]
\end{enumerate}

特に $L^p$ 空間では，
Hölder の不等式と Minkowski の不等式が
積や和の計算を支え，
Riesz--Fischer の定理が完備性を保証した．
次章では $p=2$ の場合をさらに詳しく見て，
角度や直交が使えるHilbert空間へ進む．
